\section{Optimal Offline Analysis}
\label{sec:offline}

To understand the theoretical limits of cost savings, we analyze the offline setting where the entire request sequence $\mathcal{R} = \{r_1, \dots, r_n\}$ is known in advance. We present two stages of analysis: Stage 1 considers only the economic quota constraint (a relaxation), while Stage 2 incorporates the physical concurrency constraint (the exact problem).

\subsection{Stage 1: The Economic Lower Bound (Quota-Only)}
\label{sec:offline-stage1}

We first consider a relaxed problem where only the fixed-window quota constraint exists (i.e., $C \to \infty$). This corresponds to the \textit{Economic Lower Bound}, as it ignores physical concurrency limitations.

\textbf{Problem Formulation.}
Let $x_i^Q \in \{0,1\}$ be the decision to route request $i$ to the subscription $S_Q$. The objective is to maximize the total avoided API cost (savings):
\begin{equation}
    \max \sum_{i \in \mathcal{R}} x_i^Q \cdot c_i^{\text{API}}
\end{equation}
subject to the quota constraint for each window $k$:
\begin{equation}
    \sum_{i: \operatorname{win}(a_i)=k} x_i^Q \le Q_{\text{win}}, \quad \forall k
\end{equation}

\textbf{Optimal Solution.}
Since each request consumes exactly one unit of quota, this problem decomposes into independent Knapsack problems (one per window) where all item weights are 1. The optimal strategy is simply to route the $Q_{\text{win}}$ requests with the highest API costs ($c_i^{\text{API}}$) to the subscription in each window.

\begin{theorem}
For the Quota-Only problem, the Greedy-by-Value strategy is optimal and can be computed in $\mathcal{O}(n \log n)$ time.
\end{theorem}

This oracle provides an optimistic lower bound on the API cost. Any valid solution to the full problem (Stage 2) must have a cost at least this high.

\subsection{Stage 2: The Physical Lower Bound (Quota + Concurrency)}
\label{sec:offline-stage2}

When we reintroduce the concurrency constraint $C$, the problem becomes NP-hard. It couples the knapsack-like quota constraint with a parallel machine scheduling problem.

\textbf{ILP Formulation.}
We formulate this as a time-indexed Integer Linear Program (ILP). We discretize time into slots $t \in \{0, \dots, T\}$ of width $\Delta$. Let $z_{i,t} \in \{0,1\}$ be a binary variable indicating if request $i$ starts processing at time slot $t$ on $S_C$.

The objective is to minimize total API cost:
\begin{equation}
    \min \sum_{i \in \mathcal{R}} x_i^A c_i^{\text{API}}
\end{equation}

Subject to:
\begin{align}
    & x_i^Q + x_i^C + x_i^A = 1, \quad \forall i \label{eq:assign} \\
    & \sum_{i: \operatorname{win}(a_i)=k} x_i^Q \le Q_{\text{win}}, \quad \forall k \label{eq:quota} \\
    & \sum_{t=\alpha_i}^{\beta_i} z_{i,t} = x_i^C, \quad \forall i \label{eq:consistency} \\
    & \sum_{i \in \mathcal{R}} \sum_{\tau=\max(\alpha_i, t-\hat{p}_i+1)}^{\min(\beta_i, t)} z_{i,\tau} \le C, \quad \forall t \label{eq:concurrency}
\end{align}

Constraint (\ref{eq:assign}) ensures each request is served. Constraint (\ref{eq:quota}) enforces the quota limit. Constraint (\ref{eq:consistency}) links the routing decision $x_i^C$ to the scheduling variables $z_{i,t}$. Constraint (\ref{eq:concurrency}) enforces the instantaneous concurrency limit $C$.

This ILP provides the true \textit{Physical Lower Bound} for the system. The gap between Stage 1 and Stage 2 costs represents the \textit{Cost of Concurrency}---the economic efficiency lost due to physical resource contention.
